% =============================================================================
% --- secciones/3-MotorAnalisis.tex ---
% =============================================================================

\section{Motor de Análisis}

\subsection{Snapping Lógico e Impedancia Sistémica}
\label{subsubsec:motor-analisis}

El análisis de la movilidad urbana requiere una infraestructura de datos que trascienda 
la representación geométrica simple. El \textbf{Motor de Análisis VFT} se fundamenta 
en la premisa de que la eficiencia de un sistema no depende de la 
\termino{distancia geodésica (línea recta)}, sino de la \textbf{conectividad operativa real} 
y los factores de fricción que condicionan el tiempo de traslado de los usuarios.

\subsection{Snapping Lógico: Resolución de la Inflación Topológica}

Un error común en la modelación de redes urbanas es la utilización de grafos teóricos 
ruidosos que no representan transferencias factibles. El Modelo VFT identifica este 
fenómeno como \textbf{inflación topológica}, donde la superposición de capas de información 
geoespacial genera nodos redundantes que carecen de valor operativo 
\cite{Garibaldi2024VFT, IMCO2019Doc}.

Para mitigar este ruido, se implementó un proceso de \textbf{Snapping Lógico} mediante el 
algoritmo \texttt{scipy.spatial.KDTree}. Este procedimiento permitió transitar de un grafo 
base desestructurado de aproximadamente 49,000 nodos a una red consolidada de 
\textbf{11,115 nodos operativos} y 45,679 aristas \cite{mainPGaribaldi2024VFTDF}. 
De acuerdo con los registros de ejecución del modelo (\textit{Reporte 00}), 
se logró una \textbf{reducción del 90\% de la inflación topológica}, garantizando que los 
cálculos de accesibilidad se realicen sobre puntos de transbordo y estaciones reales. 
Este rigor metodológico alinea la investigación con las tendencias actuales de Ciencia 
de Datos aplicada al urbanismo.

\subsection{Impedancia Sistémica y la Ecuación VFT}
\label{subsubsec:impendancia-VFT}

Una vez consolidada la topología de la red, el motor evalúa la \textbf{impedancia sistémica}, 
definida como el conjunto de factores físicos y operativos que restringen el derecho humano 
a la movilidad. El cambio de paradigma propuesto sustituye el enfoque de 
\termino{flujo vehicular} por uno de \termino{interacción de personas}, 
priorizando la eficiencia del sistema sobre la velocidad del objeto individual.

La piedra angular de este análisis es la \textbf{Ecuación VFT}, que cuantifica el tiempo 
de viaje efectivo ($T_t$) considerando la degradación por fricción operativa:

\begin{equation}
    T_t = \left( \frac{D}{V} \right) \times C_f
\end{equation}

Donde $D$ es la distancia, $V$ la velocidad teórica y 
\textbf{$C_f$ el Coeficiente de Fricción} \cite{Garibaldi2024VFT}. El valor de $C_f$ 
es la variable crítica que explica por qué los sistemas de superficie no confinados 
(como el transporte concesionado) colapsan ante la demanda, mientras que el 
\textbf{confinamiento masivo} se presenta como la única vía técnica para reducir 
la impedancia y garantizar tiempos de traslado competitivos.

\subsection{Comparativa de Eficiencia y Capacidad Vial}
\label{subsubsec:comparativa-eficiencia}

La fundamentación técnica del motor se refuerza al contrastar la capacidad de los modos 
de transporte. Mientras que el automóvil particular presenta una eficiencia marginal de 
\textbf{2,000 pers/hr} por carril (con una ocupación promedio de 1.2 personas), los 
sistemas masivos y el transporte público pueden movilizar entre 
\textbf{19,000 y 43,000 pers/hr} \cite{IMCO2019Doc}. 

Los datos del IMCO subrayan la ineficiencia del modelo actual: el 68.1\% de los viajes 
en auto en la ZMVM transportan un solo pasajero, ocupando un espacio vial desproporcionado 
que genera externalidades equivalentes al \textbf{4\% del PIB nacional} \cite{IMCO2019Doc}. 
El Motor VFT permite mapear estas ineficiencias, identificando que el costo del transporte 
consume hasta el \textbf{19\% del ingreso} en el cuarto contorno, lo que fundamenta la 
urgencia de una reconfiguración hacia sistemas de alta capacidad y baja fricción.
